Training NeRFs to reach photorealistic scene quality can take several days on professional high-end GPU hardware, as is the case with e.g. Mip-NeRF, a NeRF-based system reaching highest PSNR\footnote{Peak signal-to-noise ratio, used to evaluate quality of reconstruction as perceived by a human} scores among NeRFs.
Also, during novel-view synthesis, ray marching is very time-consuming as the neural network has to be queried very often and the algorithm cannot determine if a ray has reached sufficient opacity. In other words, a ray always has to travel through the entire implicit radiance field representation of the neural network.
This has led to some implementations of NeRFs switching to explicit radiance field representations and losing the neural network for novel-view synthesis allthogether (e.g. ''Plenoxels'') which comes at the inherent cost of reduced image quality due to fewer and less precise points in the radiance field.
These approaches do allow training times of 20-30 minutes and make it possible to come close to real-time novel-view synthesis reaching around 10-15 FPS at 1080p resolution. 
Other downsides of NeRFs include possibly having to sample empty space along the ray, as well as not being able to easily edit or amend a reconstructed scene as these details are all hidden in the weights of the neural network (''black box''). 
So either NeRFs allow state-of-the-art reconstruction quality at the cost of long training and synthesis times while only being able to synthesize single images or they have to reduce reconstruction quality to be able to reduce training times and come somewhat close to real-time novel-view synthesis.

To overcome the mentioned drawbacks of NeRFs, Kerbl et al \cite{kerbl3DGaussianSplatting2023} at INRIA in France introduced 3D Gaussian Splatting (3DGS) in late 2023, which makes it a very recent technology. Yet, their paper has had fundamental impact in the field of 3D reconstruction and elsewhere standing at nearly 11.000 citations \cite{researchGateKerbl} in less than three years. 

Kerbl et al managed to to combine previously existing concepts of splatting and 3D Gaussians to create a learnable explicit radiance field representation. Key to this approach is the fact that 3D Gaussians are differentiable and can be used in standard backpropagation procedures using gradient descent. Most importantly, this means that 3DGS does not use neural networks at all.

Using 3D Gaussians, a concept which we will explain in the next chapter, as the base primitive for 3D reconstruction furthermore allows for arbitrary precision and does not suffer from issues explicit radiance field representation had when used in NeRFs. Using 3D Gaussians also enables direct editability of the scene, such as additions or lighting changes making it much more accessible and subsequently usable for downstream applications.

The other important cornerstone of Kerbl et al's publication was the major engineering achievement of mapping this new idea perfectly onto modern GPU hardware for both, 3D reconstruction and novel-view synthesis, resulting in speed-ups that had not been seen before in this field. The centerpiece here is the tile-based rasterizer, which we will also cover in-depth in chapter \ref{ch:rasterizer}.

This way they created the first explicit, differentiable, unstructured and primitive-based radiance field architecture finally allowing real-time novel-view synthesis at 1080p resolution while matching or even improving scene quality in comparison to state-of-the-art NeRF architectures.

This paper will explain the underlying mechanics and concepts of 3DGS in detail. The original paper by Kerbl et al appears very packed leaving out many details, moving them to a project website and a GitHub repository.
Even many of the subsequent papers do not discuss the algorithmic approaches - which in fact have only been minorly altered since inception - so that this seminar paper takes a deep dive for an audience outside of computer graphics which is currently not being taught at Fernuniversität Hagen while emphasizing the machine learning/AI aspects.
We will begin with 3D Gaussians as the fundamental data structure and continue with the training and novel-view synthesis processes following the original paper closely before showing selected improvements to the algorithm and giving insights into 3DGS applications with a focus on dynamic scene reconstruction.